%%%%%%%%%%  scan idx 47  |  folio 36 (continued)  %%%%%%%%%%

\underline{Definition 1.4.} (1) \underline{Let $R$ be an object of
$\underline C_\Lambda$. A small extension of $R$ in
$\underline C_\Lambda$ is a surjective map
$R'\xrightarrow{f}R$ in $\underline C_\Lambda$ such that
$(\ker f)^2=0$ and $\ker f\simeq k$. In other words,
$R'\xrightarrow{f}R$ is a small extension of $R$ if}
\[
0\longrightarrow k\xrightarrow{t}R'\xrightarrow{f}R
 \longrightarrow0
\]
\underline{is exact for some $t\in R'$ with $t^2=0$.}

(2) \underline{Let $\underline A$ be a cofibered groupoid over
$\underline C_\Lambda$, and let $a$ be an object of $\underline A$.
A small prolongation of $a$ is an arrow
$a'\xrightarrow{\alpha}a$ in $\underline A$ such that $P(\alpha)$ is
a small extension in $\underline C_\Lambda$. An arrow
$a'\xrightarrow{\alpha}a$ will be called versal for small
prolongations of $a$ if, for every small prolongation
$a'_1\longrightarrow a$, there exists a commutative diagram}
\[
\begin{tikzcd}[row sep=large,column sep=large]
a' \arrow[rr] \arrow[dr,"\alpha"'] & & a'_1 \arrow[dl] \\
& a &
\end{tikzcd}
\]

\underline{Proposition 1.5. Let $\underline A$ be a semi-homogeneous
cofibered groupoid over $\underline C_\Lambda$.}

(1) \underline{Let}
\[
\begin{tikzcd}[column sep=large]
& R \\
R\times_k k[\varepsilon]
 \arrow[ur,"\pi_1"] \arrow[dr,"\pi_2"'] & \\
& k[\varepsilon]
\end{tikzcd}
\]
\underline{be the projection maps in $\underline C_\Lambda$, and let
$a'\xrightarrow{\alpha}a$ be a $\pi_1$-map. Then there exists a map
$a\xrightarrow{\beta}a'$ such that $\alpha\beta=1_a$ if and only if
there exists a map $a\longrightarrow(\pi_2)_*a'$.}

%%%%%%%%%%  scan idx 48  |  folio 37  %%%%%%%%%%

(2) \underline{Let $R'\xrightarrow{f}R$ be a small extension in
$\underline C_\Lambda$, and consider a commutative diagram}
\[
\begin{tikzcd}[row sep=large,column sep=small]
& a'' \arrow[dl,"\alpha'"'] \arrow[dr,"\alpha'_1"]
  & & & & & R'\times_RR' \arrow[dl] \arrow[dr] & \\
a' \arrow[dr,"\alpha"'] & & a' \arrow[dl,"\alpha_1"]
  & & \underline{\text{above}} & R' \arrow[dr] & & R' \arrow[dl] \\
& a & & & & & R &
\end{tikzcd}
\]
\underline{Define the canonical map
$\widetilde f:R'\times_RR'\longrightarrow k[\ker f]$ by}
\[
 \widetilde f(r'_1,r'_2)=r'_1(0)+(r'_1-r'_2).
\]
\underline{If there exists a map
$a'\longrightarrow(\widetilde f)_*(a'')$, then there exists a map
$\gamma:a'\longrightarrow a'$ such that
$\alpha=\alpha_1\gamma$.}

\underline{Proof.} (1) The ``only if'' part is clear. Assume conversely
that there exists a map
$\varphi:a\longrightarrow(\pi_2)_*a'$. Set
\[
 h:R\longrightarrow R\times_k k[\varepsilon],
 \qquad h=1\times P(\varphi).
\]
The composite $\pi_1h$ is the identity. We therefore obtain maps
\[
 a\mathrel{\mathop{\rightleftarrows}^{\gamma}_{\alpha_1}}h_*a,
 \qquad \alpha_1\gamma=1_a,
 \qquad P(\gamma)=h,
 \qquad P(\alpha_1)=\pi_1.
\]
On the other hand, $\pi_2h=P(\varphi)$, so there exists a
$\pi_2$-map
\[
 h_*a\longrightarrow(\pi_2)_*a'.
\]
Thus $(\pi_2)_*(h_*a)=(\pi_2)_*a'$, and 1.2(2) gives a
$\underline C_\Lambda$-map $u:h_*a\longrightarrow a'$ such that
$\alpha u=\alpha_1$. Define $\beta=u\gamma$. Then
\[
 \alpha\beta=\alpha u\gamma=\alpha_1\gamma=1_a.
\]

(2) The map
\[
 1\times\widetilde f:R'\times_RR'
 \longrightarrow R'\times_k k[\ker f]
\]
is an isomorphism compatible with the projections onto the first
factor, and $\pi_2(1\times\widetilde f)=\widetilde f$, where
$\pi_2:R'\times_k k[\ker f]\longrightarrow k[\ker f]$ is the second
projection. Since $R'\xrightarrow fR$ is a small extension,
$k[\ker f]\simeq k[\varepsilon]$. Hence, if there exists a map
\[
 a'\longrightarrow(\widetilde f)_*a''
   =(\pi_2)_*\bigl((1\times\widetilde f)_*a''\bigr),
\]
part (1) yields a map $\gamma':a'\longrightarrow a''$ such that
$\alpha'\gamma'=1_{a'}$. Define
$\gamma=\alpha'_1\gamma':a'\longrightarrow a'$. Then
$\alpha_1\gamma=\alpha$.

%%%%%%%%%%  scan idx 49  |  folio 38  %%%%%%%%%%

\underline{Corollary 1.6. Let $\underline A$ be a semi-homogeneous
cofibered groupoid over $\underline C_\Lambda$, and consider a diagram}
\[
\begin{tikzcd}[column sep=large]
a' \arrow[dr,"\alpha"'] & & a_1 \arrow[dl,"\alpha_1"] \\
& a &
\end{tikzcd}
\]
\underline{in $\underline A$, where $\alpha_1$ is a small
prolongation. Assume that $a$ is versal for small prolongations of
$e$. Then there exists a map $\rho:a'\longrightarrow a_1$ such that
$\alpha=\alpha_1\rho$ if and only if there exists a map
$f:P(a')\longrightarrow P(a_1)$ such that
$P(\alpha)=P(\alpha_1)f$ in $\underline C_\Lambda$.}

\underline{Proof.} Assume that $P(\alpha)=P(\alpha_1)f$, where
$f:P(a')\longrightarrow P(a_1)$. Replacing $a'$ by its direct image
under $f$, we may and shall assume that
$P(\alpha)=P(\alpha_1)$. Put $R'=P(a_1)$ and $R=P(a)$. Since
$\underline A$ is semi-homogeneous, we obtain a commutative diagram
\[
\begin{tikzcd}[row sep=large,column sep=small]
& a'' \arrow[dl,"\beta"'] \arrow[dr,"\beta_1"]
  & & & & & R'\times_RR' \arrow[dl] \arrow[dr] & \\
a' \arrow[dr,"\alpha"'] & & a_1 \arrow[dl,"\alpha_1"]
  & & \text{above} & R' \arrow[dr] & & R' \arrow[dl] \\
& a & & & & & R &
\end{tikzcd}
\]
Since $a$ is versal for small prolongations of $e$, so is $a'$;
the assertion therefore follows from 1.5(2).

